\subsubsection{Realized thinning and selected parameters}
Table~\ref{tab:s04b-thinning} counts actual rejected proposals over each complete generated history, including warmup. The narrow/short immigrant setting changes almost none of the realized event paths: nineteen histories have no rejected proposal and the remaining history rejects one immigrant. The separate random streams in the supplied generator imply that each zero-rejection immigrant-only path is identical to its same-seed uninhibited counterpart under the same runtime and other settings. This is a source-level coupling argument; no counterfactual generator replay is claimed. It does not make the two generating laws equal. The narrow/short total setting rejects offspring in every history and is a different control.

\begin{table}[htbp]\centering\small
\caption{S04B realized thinning. Rejected/proposed totals aggregate the whole histories in each regime. Zero rej.\ and Zero off.\ count histories with no rejected proposals and no accepted offspring, respectively. Mean counts use only the stated training and test windows. S4/S5 are narrow/short and broad/long immigrant attenuation; S6/S7 are the corresponding total-attenuation laws.}\label{tab:s04b-thinning}
\begin{tabular}{@{}lrrrrrr@{}}\toprule
Regime & Imm. rejected & Off. rejected & Zero rej. & Zero off. & Train $N$ & Test $N$\\\midrule
S4 & 1/2967 & 0/1120 & 19 & 0 & 86.65 & 109.40\\
S5 & 2518/3018 & 0/202 & 0 & 0 & 15.15 & 16.90\\
S6 & 1/3102 & 134/1335 & 0 & 0 & 95.80 & 111.55\\
S7 & 2525/3093 & 165/168 & 0 & 17 & 13.30 & 13.25\\
\bottomrule\end{tabular}\end{table}

Broad/long thinning also changes the overall rate strongly. These are sparse-event controls against a floor whose nuisance backbone was fixed, not evidence that inhibition beats a trained uninhibited procedure. The separate background-intercept sensitivity below measures that comparator dependence.

Table~\ref{tab:s04b-parameters} reports the selected raw fits before the complexity gate. The public package supplies every selected vector and solver trace. Boundary concentration is visible directly; an active fit is not parameter recovery. At zero strength, range and memory disappear from the likelihood, so their displayed values are arbitrary representatives of the same null intensity.
\begin{table}[H]\centering\small
\caption{S04B selected raw parameters: median above [minimum, maximum]. $I_+$ and $T_+$ denote LI+ immigrant and LI+ total. Each row uses 20 histories. $B$ counts any box boundary, including zero strength; $Z$ counts zero-strength fits. Boundary counts use the native coordinates $(\eta,\log\ell,\log\psi)$; exact equality and the recorded numerical tolerance give identical counts.}\label{tab:s04b-parameters}
\setlength{\tabcolsep}{4pt}
\begin{tabular}{@{}llcccrr@{}}\toprule
Regime & Model & $\hat\eta$ & $\hat\ell$ (km) & $\hat\psi$ (days) & $B$ & $Z$\\\midrule
S4 & LI & \shortstack{4.00\\{[}0.00, 4.00{]}} & \shortstack{57.76\\{[}5.00, 150.00{]}} & \shortstack{1.329\\{[}0.042, 365.000{]}} & 18 & 1\\
 & $I_+$ & \shortstack{2.27\\{[}0.02, 4.00{]}} & \shortstack{19.15\\{[}5.00, 150.00{]}} & \shortstack{1.645\\{[}0.042, 365.000{]}} & 14 & 0\\
 & $T_+$ & \shortstack{0.16\\{[}0.00, 0.89{]}} & \shortstack{5.62\\{[}5.00, 150.00{]}} & \shortstack{0.363\\{[}0.042, 365.000{]}} & 16 & 1\\
\addlinespace
S5 & LI & \shortstack{2.71\\{[}0.78, 4.00{]}} & \shortstack{67.19\\{[}36.45, 150.00{]}} & \shortstack{46.791\\{[}30.108, 115.698{]}} & 5 & 0\\
 & $I_+$ & \shortstack{2.00\\{[}0.56, 4.00{]}} & \shortstack{54.69\\{[}31.15, 150.00{]}} & \shortstack{31.440\\{[}19.086, 61.939{]}} & 4 & 0\\
 & $T_+$ & \shortstack{0.23\\{[}0.00, 4.00{]}} & \shortstack{95.27\\{[}5.00, 150.00{]}} & \shortstack{139.075\\{[}7.000, 365.000{]}} & 12 & 1\\
\addlinespace
S6 & LI & \shortstack{4.00\\{[}0.01, 4.00{]}} & \shortstack{28.58\\{[}5.00, 150.00{]}} & \shortstack{0.605\\{[}0.042, 365.000{]}} & 18 & 0\\
 & $I_+$ & \shortstack{2.43\\{[}0.03, 4.00{]}} & \shortstack{49.82\\{[}5.00, 150.00{]}} & \shortstack{0.571\\{[}0.042, 24.086{]}} & 16 & 0\\
 & $T_+$ & \shortstack{0.23\\{[}0.01, 2.70{]}} & \shortstack{5.28\\{[}5.00, 150.00{]}} & \shortstack{0.111\\{[}0.042, 34.715{]}} & 12 & 0\\
\addlinespace
S7 & LI & \shortstack{2.98\\{[}0.86, 4.00{]}} & \shortstack{59.83\\{[}28.30, 150.00{]}} & \shortstack{42.448\\{[}14.873, 155.075{]}} & 9 & 0\\
 & $I_+$ & \shortstack{2.19\\{[}0.51, 4.00{]}} & \shortstack{47.43\\{[}26.33, 99.70{]}} & \shortstack{28.309\\{[}11.353, 170.699{]}} & 5 & 0\\
 & $T_+$ & \shortstack{2.14\\{[}0.49, 4.00{]}} & \shortstack{49.92\\{[}26.86, 117.12{]}} & \shortstack{28.705\\{[}11.164, 182.512{]}} & 4 & 0\\
\bottomrule\end{tabular}\end{table}

The saved narrow/short immigrant example with unit-mark LI illustrates start dependence. A scale-matched starting point reaches zero strength with a positive one-sided negative-log-likelihood derivative in strength and zero shape derivatives; those signs satisfy the local constrained optimum conditions. Another prescribed start attains a lower training objective and is selected. Thus the null-face termination is not by itself a solver defect or a global-optimality certificate. The candidate also uses unit stress marks, so matching its scales to a magnitude-scaled generator does not make it the true generating law.
